% BAN TIENG VIET cua 03_guarantee.tex — viet lai 2026-10-12 theo academic-prose.
% Moi con so, cite, label, ref, equation giu NGUYEN so voi ban truoc.
% Thuat ngu theo glossary.tex. KHONG them/bot claim.
\section{Khi nào được kết tội một nông dân?}
\label{sec:guarantee}

Hai mục trước cho thấy lớp kiểm toán vật lý vừa bỏ lọt bên gian lận tinh vi vừa có nguy
cơ kết tội oan người trung thực. Mục này xử lý vế thứ hai trước khi tổng hợp thành khẳng định điểm mù không thể tránh.

\subsection{Độ sâu tưới chấp nhận được là một dải, không phải một con số}
\label{sec:umax}

Nếu chỉ lấy giới hạn từ một hồ sơ đất và một điều kiện đầu duy nhất thì cách
tính cho ra đúng một con số, $u\le 61$~mm, là khai báo lớn nhất không bao giờ
gây báo động. Vì $V_4$ kích hoạt khi $u_i>c_{t_i}+\varepsilon$ và
$c_t=d_t\cdot\mathrm{span}$, giới hạn tại một ruộng có độ cạn kiệt ban đầu
$w_0$ và tưới không muộn hơn $d^{*}$ là
\begin{equation}
u_{\max}=\max(d^{*},w_0)\cdot\mathrm{span}+\varepsilon .
\label{eq:umax}
\end{equation}
trong đó:
\begin{symbolist}
\item[$u_{\max}$] độ sâu khai báo lớn nhất, tính bằng mm, mà bộ lọc không bao giờ đánh dấu;
\item[$d^{*}$] độ cạn kiệt tại giờ tưới muộn nhất được xét, là phân số không thứ nguyên;
\item[$w_0$] độ cạn kiệt ban đầu của ruộng, cùng đơn vị phân số;
\item[$\mathrm{span}$] biên độ của thang cạn kiệt, bằng $FC-WP$ tính bằng mm, đã định nghĩa ở Mục~\ref{sec:model};
\item[$\varepsilon$] dung sai công bố cho sai số mô hình, bằng $5$~mm như ở Mục~\ref{sec:filters}.
\end{symbolist}

Giới hạn này được viết theo tỷ lệ cạn kiệt quy về tầng rễ chứ không theo một
cao độ tuyệt đối, nên khác với các ngưỡng mực nước tuyệt đối đang dùng trong thực
hành, nó không trôi khi mặt ruộng sụt lún; mệnh đề và độ dịch đo được cho đồng bằng
ở Phụ lục~S4.

Từ đây rút ra hai điều chỉnh, cả hai đều được kiểm chứng bằng cách quét toán tử
phát lại chứ không chỉ bằng đại số (dẫn xuất ở Phụ lục~S14). Thứ nhất,
\eqref{eq:umax} dùng $\max$ chứ không dùng hiệu, nên một ruộng khô ban đầu lại
chấp nhận khai báo \emph{lớn hơn} một ruộng ướt, còn biểu thức dạng hiệu thì
không tương ứng với trạng thái vật lý nào. Thứ hai, giới hạn này chỉ chặn một
\emph{lần tưới đơn lẻ} chứ không chặn cả vụ, nên một bên gian lận phát nhiều khai báo
dưới ngưỡng có thể thổi phồng tổng lượng mà không bị giới hạn nào chặn lại.
Bảng~S2 của Phụ lục đưa ra dải cần được báo cáo thay thế, $[49{,}0; 86{,}0]$~mm
trên toàn $\Theta$.

\subsection{Chỉ kết tội khi mọi kịch bản mô hình đều cho thấy vi phạm}
\label{sec:robust}

Kiểm toán viên không biết đất của nông dân, không biết giai đoạn sinh trưởng và
cũng không biết hiệu suất của hệ thống dẫn nước. Gọi $\Theta$ là tập bất định và
$V_\theta(L)$ là tập vi phạm dưới cấu hình $\theta$. Bộ phát hiện mô hình điểm
chỉ đánh giá một $\theta_0$ mặc định rồi coi mọi vi phạm là có tội. Chúng tôi
thay nó bằng một phép giao:
\begin{equation}
\begin{aligned}
V_{\mathrm{rob}}(L)&=\bigcap_{\theta\in\Theta}V_\theta(L),\\
\Theta&=\mathcal{S}\times\mathcal{K}\times\mathcal{E}\times\mathcal{W},
\end{aligned}
\label{eq:robust}
\end{equation}
trong đó:
\begin{symbolist}
\item[$V_{\mathrm{rob}}(L)$] tập vi phạm bền vững của nhật ký $L$, chỉ khác rỗng khi $L$ vi phạm dưới mọi cấu hình;
\item[$\Theta$] tập bất định, là tích của bốn tập thành phần dưới đây;
\item[$\mathcal{S}$] tập loại đất, gồm $5$ loại;
\item[$\mathcal{K}$] tập chế độ hệ số cây trồng, gồm $2$ chế độ;
\item[$\mathcal{E}$] tập mức hiệu suất tưới, gồm $3$ mức;
\item[$\mathcal{W}$] tập trạng thái đầu, gồm $5$ trạng thái.
\end{symbolist}

Do đó $|\Theta|=150$ và mỗi nhật ký tốn $150$ lần phát lại. Ngữ nghĩa của phán quyết được chọn cố ý lệch về phía bảo vệ nông dân, tức ưu tiên
tránh kết tội oan hơn là tối đa năng lực phát hiện:
\begin{equation}
a(L)=
\begin{cases}
\text{kết tội}, & V_{\mathrm{rob}}(L)\neq\emptyset,\\
\text{chờ xác minh}, & \text{trường hợp còn lại}.
\end{cases}
\label{eq:verdict}
\end{equation}
trong đó $a(L)$ là phán quyết dành cho nhật ký $L$. Phán quyết ``chờ xác minh''
nghĩa là phải xác minh thêm tại thực địa, chứ hệ thống không tự kết luận gian lận.

\begin{proposition}[Không kết tội oan trong $\Theta$]
\label{prop:noalarm}
Nếu cấu hình thật của ruộng $\theta^{*}$ thuộc $\Theta$ và $L$ tuân thủ dưới
$\theta^{*}$, thì $V_{\mathrm{rob}}(L)=\emptyset$ và nông dân không bị kết tội.
\end{proposition}

\begin{proof}
$V_{\mathrm{rob}}(L)\subseteq V_{\theta^{*}}(L)=\emptyset$.
\end{proof}

Mệnh đề~\ref{prop:noalarm} là cơ sở để lớp này triển khai được trong thực tế. Sai
số mô hình có thể làm giảm năng lực phát hiện nhưng không thể tạo ra một lời kết
tội oan, và cái giá phải trả cho đảm bảo đó được nêu rõ ở phần sau. Hai chi tiết cài đặt có thể âm thầm phá vỡ mệnh đề nếu làm sai, đó là cách xử lý
những cấu hình trong tập bất định mà tại đó không bộ lọc nào kích hoạt, và độ trải
của tập trạng thái đầu $\mathcal{W}$ (Phụ lục~S7).

\subsection{Ngưỡng chênh lệch bơm tỉ lệ thuận}

Quy tắc tuyệt đối $G>\tau_0$ với $\tau_0=5$~mm ngầm giả định nhật ký là bản sao
trung thực của lượng nước đã bơm. Giả định này không đúng, vì quên ghi, làm tròn
đến $5$~mm gần nhất và ghi trễ một ngày đều làm tăng $G$ mà không hề có hành vi
sai trái nào. Vì vậy ngưỡng phải co giãn theo số lần tưới được ghi:
\begin{equation}
\begin{aligned}
\tau(n)&=\delta_{\mathrm{rel}}M+z\sqrt{n}\,\sigma_{\mathrm{event}},\\
\sigma_{\mathrm{event}}^{2}&=p(1-p)\bar u^{2}+\tfrac{r^{2}}{12}
+q(1-q)(s\bar u)^{2},
\end{aligned}
\label{eq:tau}
\end{equation}
trong đó:
\begin{symbolist}
\item[$\tau(n)$] ngưỡng chênh lệch bơm cho phép, tính bằng mm, khi nhật ký có $n$ lần tưới;
\item[$n$] số lần tưới được ghi trong nhật ký;
\item[$\delta_{\mathrm{rel}}$] tỷ lệ sai lệch tương đối cho phép trên lượng bơm thật $M$, bằng $2\%$;
\item[$z$] hệ số nhân của độ lệch chuẩn, đặt bằng $3$;
\item[$\sigma_{\mathrm{event}}$] độ lệch chuẩn của sai số ghi chép trên mỗi lần tưới, tính bằng mm;
\item[$p$] tỷ lệ quên ghi một lần tưới;
\item[$r$] bước làm tròn độ sâu, bằng $5$~mm;
\item[$q$] tỷ lệ ghi sai độ sâu, và $s$ là biên độ tương đối của lỗi đó;
\item[$\bar u$] độ sâu điển hình của một lần tưới, bằng $30$~mm.
\end{symbolist} Hai quy tắc không phải hai trường hợp của cùng một công thức, và
không nên đọc \eqref{eq:tau} như dạng tổng quát của $\tau_0$. Khi
$n\to0$ và $\sigma_{\mathrm{event}}\to0$, biểu thức trên tiến về
$\delta_{\mathrm{rel}}M$, tức $2\%$ lượng bơm thật, chứ không tiến về $5$~mm; hai
ngưỡng chỉ trùng nhau nếu $M$ đúng bằng $250$~mm. Chúng khác nhau về bản chất:
$\tau_0=5$~mm là hằng số tuyệt đối, chỉ đúng khi nhật ký được giả định là bản sao
chính xác của số bơm, còn $\tau(n)$ co giãn theo quy mô $M$ và theo số sự kiện
$n$. Trong thực nghiệm chúng tôi dùng $\tau(n)$ cho mọi phán quyết và giữ
$\tau_0$ làm mốc so sánh, để chỉ ra cái giá của giả định bản sao trung thực.

\begin{remark}[Bộ lọc thứ ba không bao giờ kích hoạt được]
\label{sec:v2reach}
Lớp vật lý có ba bộ lọc, nhưng bộ thứ ba --- đánh dấu ruộng cạn kiệt xuống dưới
điểm héo --- không bao giờ kích hoạt được, vì hai lý do độc lập. Lý do thứ nhất
thuộc về cách cài đặt: ngưỡng của nó là $1{,}02$, trong khi chỉ số được báo cáo đã
bị cắt ngọn về một đoạn cố định, nên điều kiện không thể nào đúng và bộ lọc có $0$
lần kích hoạt trong cả $281$ nhật ký. Lý do thứ hai thuộc về thực tế đồng ruộng:
chúng tôi sửa lỗi đó bằng cách dùng độ cạn kiệt chưa cắt ngọn, rồi đo lại trên
$150$ cấu hình với nhật ký rỗng, tức điều kiện khô hạn nhất có thể, thì giá trị này
cũng chỉ đạt đỉnh $0{,}9987$. Nghĩa là một ruộng lúa vụ khô ở đồng bằng không cạn
tới mức chạm điểm héo (Bảng~S3 của Phụ lục, dẫn xuất ở S13). Bộ lọc này vì vậy chỉ
là một chốt an toàn chống trạng thái vô lý, không phải bằng chứng về năng lực phát
hiện, và lớp này chỉ có hai bộ lọc vật lý thực sự hoạt động.
\end{remark}
